\subsection{Reference-line checks and review exclusions}
\label{app:ccl_detailed_sensitivities}

Human review checked whether the accepted lines state the challenged claim. The initial 40-input check approved 35 line selections; after corrections, a fresh, disjoint sample approved 38 of 40 and met the review criterion. Appendix~\ref{app:recent-human-review} gives the full review process and dispositions.

The current 100-input primary set and 131-input scorable set exclude every record with a retained negative or uncertain source, witness, or formalization review, and all three records with adverse judgments on their current line selections. Reapplying these exclusion lists leaves the current sets and scores unchanged. Appendix~\ref{app:recent-human-review} records the individual decisions; \path{bugbank/reviews/version_resolution/README.md} supplies the complete exclusion lists and saved-answer sensitivity. Reference-line revisions and their replay are documented in \path{benchmark/README.md}.

\subsection{Examining first-choice misses}
\label{app:ccl_miss_diagnostic}

A random sample of 30 was drawn from Sonnet 5's 58 readable first-choice misses in the earlier 191-input evaluation, using seed 20260921. LLM reviewers, with access to the source text and reference labels, classified the selected claims and assessed the accompanying explanations.

The final review identified 12 selections of another false assertion, 15 lines that were not false as written, and three source or annotation problems. The first LLM reviewer approved the supplied explanation in 11 of the 12 alternative-assertion cases. The three source or annotation problems were two omitted restatements of the challenged proposition and one garbled source.

The released records preserve the original classifications, later adjudications, sampling intervals, and disagreements on M08, M13, and M20. This analysis retained the original benchmark scores. The sampled responses and scientific reviews are in \path{bugbank/reviews/sonnet_misses/}; \path{bugbank/reviews/README.md} explains their numerical reproduction.

